\chapter{Radiation Dose and Safety Comparison Chart}
\label{app:radiation-chart}

\noindent
One of the most common questions patients ask is: \emph{``Does this scan use radiation, and how much?''} Radiation dose is measured in \textbf{millisieverts (mSv)}. Every person naturally receives about \textbf{3~mSv per year} (roughly \textbf{0.008~mSv per day}) from everyday natural background sources such as cosmic rays, soil, radon, and food. For comparison, a single coast-to-coast commercial airplane flight adds about \textbf{0.03--0.04~mSv}.

\section{Modalities With Zero Ionizing Radiation (0 mSv)}
These imaging families use magnets, sound waves, or gentle light rather than X-rays or radioactive tracers. They carry \textbf{zero ionizing radiation risk}, though each has its own safety rules:

{\renewcommand{\arraystretch}{1.25}\noindent
\begin{tabularx}{\textwidth}{@{}p{0.28\textwidth} p{0.34\textwidth} X@{}}
\toprule
\textbf{Modality Family} & \textbf{Example Chapters} & \textbf{Primary Safety Focus} \\
\midrule
Magnetic Resonance (MRI) & MRI (Ch.~\ref{ch:infographic-mri}), Breast MRI (Ch.~\ref{ch:infographic-breast-mri}), Cardiac MRI (Ch.~\ref{ch:infographic-cardiac-mri}), DWI, MRA, MRV, MRCP, MRE & Strict metal and implant screening; kidney review if gadolinium dye is needed \\
Passive Brain Magnetism & MEG (Ch.~\ref{ch:infographic-meg}) & Metal removal for clean signal; completely passive \\
Sound-Wave (Ultrasound) & Ultrasound (Ch.~\ref{ch:infographic-ultrasound}), Prenatal US (Ch.~\ref{ch:infographic-obstetric-ultrasound}), Doppler, CEUS, Elastography, POCUS, Echo, EUS, B-Scan, UBM & Gentle thermal/mechanical limits; sedation precautions only for EUS \\
Optical \& Endoscopy & Endoscopy (Ch.~\ref{ch:infographic-endoscopy}), Capsule, Slit-Lamp, Fundus, FFA, ICGA, OCT, OCTA, Corneal Maps, Dermoscopy (Ch.~\ref{ch:infographic-dermoscopy}), Confocal, HSI & Sedation/bleeding checks for GI scopes; dye allergy check for FFA/ICGA; eye-safe laser limits \\
Emerging Non-Ionizing Research Tools & MPI (Ch.~\ref{ch:infographic-mpi}), Raman, OAT, PAI, EIT, Microwave, Terahertz & Investigational IRB oversight; laser eyewear or magnet/thermal limits \\
\bottomrule
\end{tabularx}}

\section{Ionizing-Radiation Modalities (Lowest to Highest Typical Dose)}
Typical adult effective doses below are approximate averages from major radiology organizations (ACR, RSNA, and ICRP). Actual doses vary with body size, modern low-dose equipment, and whether multiple phases are scanned.

\subsection{Very Low to Low Dose (Under 2 mSv)}
{\renewcommand{\arraystretch}{1.2}\noindent
\begin{tabularx}{\textwidth}{@{}p{0.31\textwidth} p{0.16\textwidth} p{0.20\textwidth} X@{}}
\toprule
\textbf{Scan Type (Chapter)} & \textbf{Typical Dose} & \textbf{Natural Background Equivalent} & \textbf{Notes / Typical Contrast} \\
\midrule
Dental X-Ray --- Bitewing (Ch.~\ref{ch:infographic-dental-radiography}) & $\sim$0.005~mSv & Less than 1~day & No contrast; ultra-low localized dose \\
Limb / Hand / Foot X-Ray (Ch.~\ref{ch:infographic-xray}) & 0.001--0.01~mSv & Less than 1--2~days & No contrast \\
DEXA Bone Density (Ch.~\ref{ch:infographic-dexa}) & 0.001--0.01~mSv & 1--2~days & No contrast; safe for routine tracking \\
Orthopantomogram / OPG (Ch.~\ref{ch:infographic-opg}) & 0.01--0.025~mSv & 1--3~days & No contrast; panoramic jaw view \\
Chest X-Ray (Ch.~\ref{ch:infographic-xray}) & 0.02--0.1~mSv & 3--10~days & No contrast; roughly 1--2 cross-country flights \\
Dual-Energy Chest X-Ray / DEI (Ch.~\ref{ch:infographic-dei}) & 0.05--0.15~mSv & 6--18~days & No contrast; rib-subtracted chest view \\
Dental / ENT CBCT (Ch.~\ref{ch:infographic-cbct}) & 0.05--0.3~mSv & 1--5~weeks & No contrast; lower dose than head CT \\
Mammography / 3D DBT (Ch.~\ref{ch:infographic-mamography}, \ref{ch:infographic-dbt}) & 0.4--0.5~mSv & About 7~weeks & No contrast; standard screening exam \\
Abdomen / Lower-Spine X-Ray (Ch.~\ref{ch:infographic-xray}) & 0.7--1.5~mSv & 3--6~months & No contrast \\
Head CT (Ch.~\ref{ch:infographic-ctscan}) & 1.5--2.0~mSv & 6--8~months & Usually no contrast for trauma/bleeding \\
\bottomrule
\end{tabularx}}

\subsection{Moderate to Higher Dose (2 mSv and Above)}
{\renewcommand{\arraystretch}{1.2}\noindent
\begin{tabularx}{\textwidth}{@{}p{0.31\textwidth} p{0.16\textwidth} p{0.20\textwidth} X@{}}
\toprule
\textbf{Scan Type (Chapter)} & \textbf{Typical Dose} & \textbf{Natural Background Equivalent} & \textbf{Notes / Typical Contrast} \\
\midrule
Cardiac CT / Calcium Score (Ch.~\ref{ch:infographic-cardiac-ct}) & 1--5~mSv & 4--20~months & None (calcium score) or IV iodine (coronary CTA) \\
Fluoroscopy / Myelography (Ch.~\ref{ch:infographic-fluroscopy}, \ref{ch:infographic-myelography}) & 2--8~mSv & 8~months--2.5~years & Barium or iodine contrast; varies with live duration \\
Nuclear Scintigraphy (Ch.~\ref{ch:infographic-nms}) & 2--6~mSv & 8~months--2~years & IV/oral radiotracer (bone, thyroid, V/Q, HIDA) \\
Chest, Belly, or Pelvis CT / DECT / CTA / CTP (Ch.~\ref{ch:infographic-ctscan}, \ref{ch:infographic-dual-energy-ct}, \ref{ch:infographic-cta}, \ref{ch:infographic-ctp}) & 5--15~mSv & 1.5--5~years & Often uses IV iodine contrast; low-dose lung CT is $\sim$1.5~mSv \\
SPECT / SPECT-CT (Ch.~\ref{ch:infographic-spect}, \ref{ch:infographic-spect-ct}) & 4--12~mSv & 1.5--4~years & IV radiotracer (+ low-dose CT for localization) \\
PET / PET-MRI (Ch.~\ref{ch:infographic-petscan}, \ref{ch:infographic-pet-mri}) & 4--7~mSv & 1.5--2.5~years & IV radiotracer only; avoids CT X-ray dose (50--80\% lower than PET/CT) \\
PET/CT Whole-Body (Ch.~\ref{ch:infographic-pet-ct}) & 7--25~mSv & 2--8~years & IV radiotracer (e.g., FDG) + CT component \\
DSA \& Complex IR Procedures (Ch.~\ref{ch:infographic-dsa}, \ref{ch:infographic-interventional-radiology}) & 5--30+~mSv & Variable (tracked live) & IV iodine contrast; high doses closely monitored for skin safety \\
\bottomrule
\end{tabularx}}

\section{Key Safety Principles to Remember}
\begin{itemize}
  \item \textbf{Benefit Always Outweighs Small Risk:} When a doctor orders an X-ray, CT, or nuclear scan for a clear medical reason, the value of finding a fracture, blood clot, stroke, or tumor far outweighs the tiny theoretical radiation risk.
  \item \textbf{ALARA (As Low As Reasonably Achievable):} Modern scanners automatically tailor radiation to your body size, use child-specific low-dose protocols, and narrow the beam strictly to the area being checked.
  \item \textbf{Pregnancy and Children:} Always tell your care team if you might be pregnant. Whenever a radiation-free ultrasound or MRI can answer the same medical question just as well, doctors choose the radiation-free test first.
  \item \textbf{Avoid Duplicate Scans:} Bring prior CDs, portal records, or outside reports to appointments so your doctor can compare old images instead of repeating a scan unnecessarily.
\end{itemize}
